\documentclass[10pt,letterpaper,twocolumn]{article}

\usepackage[utf8]{inputenc}
\usepackage[T1]{fontenc}
\usepackage{newtxtext,newtxmath}
\usepackage[scaled=0.92]{helvet}
\usepackage[top=0.76in,bottom=0.82in,left=0.72in,right=0.72in,columnsep=0.23in]{geometry}
\usepackage{microtype}
\usepackage{amsmath,mathtools}
\usepackage{graphicx}
\usepackage{booktabs}
\usepackage{longtable}
\usepackage{array}
\usepackage{enumitem}
\usepackage{titlesec}
\usepackage{xcolor}
\usepackage[hidelinks]{hyperref}
\usepackage{url}
\usepackage{fancyhdr}
\usepackage{setspace}

\setstretch{0.995}
\setlength{\parindent}{0.9em}
\setlength{\parskip}{0pt}
\setlength{\columnsep}{0.23in}
\setlist[itemize]{leftmargin=1.2em,itemsep=0.2em,topsep=0.35em}
\setlist[enumerate]{leftmargin=1.3em,itemsep=0.2em,topsep=0.35em}
\setlength{\abovedisplayskip}{5pt plus 2pt minus 2pt}
\setlength{\belowdisplayskip}{5pt plus 2pt minus 2pt}
\setlength{\abovedisplayshortskip}{4pt plus 2pt minus 2pt}
\setlength{\belowdisplayshortskip}{4pt plus 2pt minus 2pt}
\setlength{\textfloatsep}{8pt plus 2pt minus 2pt}
\setlength{\floatsep}{7pt plus 2pt minus 2pt}
\setlength{\intextsep}{7pt plus 2pt minus 2pt}
\setlength{\emergencystretch}{1.5em}
\raggedbottom

\definecolor{TitleRule}{HTML}{D9DEE7}
\definecolor{TitleNote}{HTML}{4B5563}
\definecolor{HeadingInk}{HTML}{1F2937}
\definecolor{HeadingRule}{HTML}{CBD5E1}

\renewcommand{\thesection}{\arabic{section}}
\renewcommand{\thesubsection}{\thesection.\arabic{subsection}}
\renewcommand{\thesubsubsection}{\thesubsection.\arabic{subsubsection}}
\setcounter{secnumdepth}{2}

\titleformat{\section}
  {\large\sffamily\bfseries\color{HeadingInk}}
  {\thesection}
  {0.65em}
  {}
\titleformat{name=\section,numberless}
  {\large\sffamily\bfseries\color{HeadingInk}}
  {}
  {0pt}
  {}
\titleformat{\subsection}
  {\normalsize\sffamily\bfseries\color{HeadingInk}}
  {\thesubsection}
  {0.55em}
  {}
\titleformat{\subsubsection}
  {\normalsize\sffamily\bfseries\color{HeadingInk}}
  {\thesubsubsection}
  {0.5em}
  {}
\titlespacing*{\section}{0pt}{0.95em}{0.35em}
\titlespacing*{\subsection}{0pt}{0.75em}{0.25em}
\titlespacing*{\subsubsection}{0pt}{0.55em}{0.2em}

\pagestyle{fancy}
\fancyhf{}
\fancyfoot[C]{\small\thepage}
\renewcommand{\headrulewidth}{0pt}

\providecommand{\tightlist}{%
  \setlength{\itemsep}{0pt}\setlength{\parskip}{0pt}}

\title{A Summary of the Geodesic-Coherence Refinement for Fourier-Domain Photon-Ring Inference}
\author{Jonathan R. Landers}
\date{}

\newcommand{\affiliationnote}{Independent researcher.\\
This note summarizes the arc across \\
\emph{A Fourier-Domain Analysis of BHEX for Photon-Ring Inference} and\\
\emph{Geodesic-Provenance Coherence Bounds for Fourier-Domain Photon-Ring Inference}, and is not affiliated\\
with the BHEX collaboration or mission team.}

\makeatletter
\renewcommand{\maketitle}{%
  \twocolumn[{%
    \begin{@twocolumnfalse}
      \vspace*{-1.3em}
      {\LARGE\sffamily\bfseries\raggedright \@title \par}
      \vspace{0.35em}
      {\normalsize\sffamily\raggedright \@author \par}
      \vspace{0.16em}
      {\small\color{TitleNote}\raggedright \affiliationnote \par}
      \vspace{0.65em}
      {\color{TitleRule}\hrule height 0.8pt}
      \vspace{0.8em}
    \end{@twocolumnfalse}
  }]
}
\makeatother

\begin{document}
\maketitle

The original Fourier-domain note [\hyperlink{ref-1}{1}] began with a simple model for long-baseline
black-hole interferometry,
\[
y = \alpha g + p + \varepsilon,
\]
where $g$ is a photon-ring template, $p$ is structured background emission, and
$\varepsilon$ is noise. Its main conceptual point was that \textbf{readability}
and \textbf{recoverability} are not the same thing: a direct amplitude heuristic can lose
clarity once the modulus nonlinearity distorts oscillation spacing, even though a structured
inverse problem can still recover the ring stably. The mismatch between those two regimes is
controlled by a ring-background coherence term together with nuisance size and noise.

The later geodesic-coherence note [\hyperlink{ref-2}{2}] asked what that coherence term really means
physically. Its answer was to lift the nuisance model from visibility space back to photon
initial-condition space and show that the relevant overlap can itself be bounded by the geometry of
escaping geodesic families. This summary note records that arc in one place, so that a reader who has
already seen the two notes can latch onto the intuitive flow and the final result.

\section{The original Fourier-domain story}

In the first note [\hyperlink{ref-1}{1}], the visibility data are modeled in a complex Hilbert space as
\[
y = \alpha g + p + \varepsilon.
\]
A simple ring-oriented heuristic reads off information directly from the visibility amplitude, while a
plasma-aware estimator explicitly projects away or jointly models structured nuisance. In the simplest
subspace version, the key overlap quantity is
\[
\mu(g,\mathcal P)=\frac{\|\Pi_{\mathcal P}g\|}{\|g\|},
\]
where $\Pi_{\mathcal P}$ is orthogonal projection onto a nuisance subspace $\mathcal P$.

The first lemma in that note made the mismatch story concrete: the discrepancy between a direct-readout
estimator and a plasma-aware estimator is controlled by a coherence term of the schematic form
\[
\mu(g,\mathcal P)\frac{\|p\|}{\|g\|},
\]
together with noise contributions. So the heuristic does not fail arbitrarily. Its bias grows when the
assumed ring template has significant projection into the nuisance family, or when the nuisance and noise
become large relative to the ring.

The accompanying proposition then sharpened the conceptual message. Because the modulus is nonlinear,
moderate contamination can shift amplitude extrema and distort oscillation spacing, so \emph{clean
readability} can break down even while a structured inverse problem remains locally identifiable and
stable. In other words, \textbf{the heuristic readability limit can occur before the recoverability
limit}. This was the first note's main structural contribution: it recast the BHEX long-baseline picture
as a source-separation problem in which visual cleanliness is not the same thing as mathematical
recoverability [\hyperlink{ref-3}{3}, \hyperlink{ref-4}{4}].

\section{Why provenance had to enter}

That first framework was intentionally abstract. It treated the nuisance class geometrically, but did not
ask where the nuisance actually comes from. The next step was to notice that photons are not arbitrary
signals in visibility space. They arise from initial conditions
\[
z=(p,\xi),
\]
where $p$ is an emission event and $\xi$ is an initial null direction, and then propagate along null
geodesics.

Let $E$ denote the set of escaping photon initial conditions. The crucial refinement is that the
observable set itself naturally splits as
\[
E = E_{\mathrm{ring}} \,\dot\cup\, E_{\mathrm{bg}},
\]
where $E_{\mathrm{ring}}$ consists of near-critical escaping geodesics that generate the photon-ring
signal, while $E_{\mathrm{bg}}$ consists of ordinary escaping geodesics that generate the broader
background. Captured photons do not contribute to the data, but the deeper point is that even within the
observable sector, ring and background come from distinct provenance classes.

This reinterprets the Fourier-domain decomposition. Instead of treating $p$ or $q$ as an arbitrary
structured contaminant, one writes the ring and the background as visibility-space images of transport
operators acting on initial-condition densities. If
\[
A_r=\mathcal F T_{E_{\mathrm{ring}}},
\qquad
A_b=\mathcal F T_{E_{\mathrm{bg}}},
\]
then one has
\[
g=A_r f_\theta,
\qquad
q=A_b h,
\]
for suitable emission densities $f_\theta$ and $h$. The nuisance family is therefore no longer generic;
it is constrained to the provenance class
\[
Q_{\mathrm{prov}}=\{A_b h:h\in\mathcal A\},
\]
for an admissible emission class $\mathcal A$.

\section{The main new theorem}

The real mathematical advance of the second note [\hyperlink{ref-2}{2}] was not just the observation that
$Q_{\mathrm{prov}}$ is smaller than an abstract nuisance class. It was an explicit upper bound on the
provenance-constrained coherence itself. For a ring template $g$, define
\[
\mu(g,Q_{\mathrm{prov}})
:=
\sup_{q\in Q_{\mathrm{prov}}\setminus\{0\}}
\frac{|\langle g,q\rangle|}{\|g\|\,\|q\|}.
\]
Assume the transport operators satisfy lower nondegeneracy bounds
\[
\|A_r f\|\ge \sigma_r\|f\|,
\qquad
\|A_b h\|\ge \sigma_b\|h\|,
\]
on the relevant classes.

If the operators admit a visibility kernel $K(u,z)$, define the cross-Gram kernel
\[
G(z,z'):=\langle K(\cdot,z),K(\cdot,z')\rangle_{\mathcal H}.
\]
Then the key theorem says that
\begin{equation}
\begin{aligned}
\mu(g,Q_{\mathrm{prov}})
&\le
\frac{1}{\sigma_r\sigma_b}
\biggl(
\int_{E_{\mathrm{ring}}}\!\int_{E_{\mathrm{bg}}}
|G(z,z')|^2\,d\nu(z)\,d\nu(z')
\biggr)^{1/2}.
\end{aligned}
\end{equation}

This is the centerpiece. The abstract overlap term from the original note is now lifted to photon
initial-condition space and bounded by a cross-correlation integral over escaping geodesic families.
That is what turns the provenance idea into a genuinely new theorem rather than a structural remark.

\section{The geometric separation corollary}

To turn the preceding theorem into a more geometric statement, the second note introduces a scalar
criticality index $\chi(z)$ measuring how ``ring-like'' an escaping trajectory is. The specific choice can
vary; it might be tied to winding depth, near-photon-region residence time, or another surrogate for
criticality.

Assume there is a gap $\Delta>0$ such that ring-generating and background-generating trajectories are
separated in this index:
\[
|\chi(z)-\chi(z')|\ge \Delta
\quad
\text{for } z\in E_{\mathrm{ring}},\ z'\in E_{\mathrm{bg}}.
\]
Assume also that the cross-Gram kernel decays exponentially with this separation:
\[
|G(z,z')|\le M e^{-\beta |\chi(z)-\chi(z')|}.
\]
Then the coherence itself obeys the explicit bound
\begin{equation}
\begin{aligned}
\mu(g,Q_{\mathrm{prov}})
&\le
\frac{M}{\sigma_r\sigma_b}e^{-\beta\Delta} \\
&\quad \cdot \sqrt{\nu(E_{\mathrm{ring}})\,\nu(E_{\mathrm{bg}})}.
\end{aligned}
\end{equation}

This is the genuinely geometric conclusion. It says that the key overlap term in the Fourier-domain
mismatch story becomes quantitatively small when the near-critical escaping flow and the ordinary
escaping flow are dynamically well separated. The stronger statement is not merely that the nuisance class
is smaller, but that its coherence with the ring can decay exponentially under a phase-space separation
hypothesis.

\section{What changed across the two notes}

Taken together, the two notes tell a coherent story.

The first note [\hyperlink{ref-1}{1}] established that direct amplitude readability can fail before
structured recoverability, with the mismatch controlled by coherence, nuisance magnitude, and noise.
The second note [\hyperlink{ref-2}{2}] identified a physically meaningful way to sharpen that coherence
term. By lifting the model to photon initial-condition space and splitting the escaping set into ring and
background provenance classes, it proved that the relevant overlap is bounded by a cross-Gram integral over
escaping geodesic families. Under a criticality-gap assumption, that overlap decays exponentially.

So the original question was:
\[
\text{Can the ring remain recoverable after it stops being cleanly readable?}
\]
The refined answer is:
\[
\text{Yes, and the degree of recoverability is controlled by}
\]
\[
\text{the geometric separation of escaping photon families.}
\]

That is the intuitive arc. The original note gave an abstract structural theory of mismatch and
robustness. The second note showed that the key overlap term in that theory is not merely formal: under
natural operator and separation assumptions, it is governed by null-geodesic phase-space geometry.

\section{Relation to BHEX and computational outlook}

This summary still sits squarely in the spirit of the BHEX long-baseline program
[\hyperlink{ref-3}{3}, \hyperlink{ref-4}{4}, \hyperlink{ref-5}{5}]. BHEX emphasizes that sufficiently
long baselines should reveal a more universal photon-ring signal beneath broader astrophysical emission.
The first note recast that intuition as a Fourier-domain structured source-separation problem. The second
note then identified a route by which the key overlap term in that source-separation problem could, in
principle, be estimated from geodesic geometry.

The framework remains abstract until one chooses an admissible emission class, a concrete criticality
index, and a practical way to estimate the cross-Gram kernel from ray tracing or simulation. But this is
exactly the bridge the two notes were moving toward: from abstract mismatch and recoverability statements
to physically informed BHEX-style inference.

\section{Key interesting points}

\begin{itemize}
\tightlist
\item The first note introduced a Fourier-domain structured source-separation framework and showed that
readability can fail before recoverability [\hyperlink{ref-1}{1}].
\item The second note lifted the nuisance model to photon initial-condition space and split the observable
sector into ring-generating and background-generating escaping geodesics [\hyperlink{ref-2}{2}].
\item The main new theorem bounded provenance-constrained coherence by a cross-Gram integral over escaping
geodesic families.
\item Under a criticality-gap assumption, that coherence decays exponentially with phase-space separation.
\item The central advance is therefore not just ``less nuisance,'' but a concrete mechanism by which
black-hole geodesic geometry sharpens the original Fourier-domain mismatch story.
\end{itemize}

\section*{References}
\addcontentsline{toc}{section}{References}

\noindent{}\hypertarget{ref-1}{}{[}1{]} Jonathan R. Landers,
\emph{A Fourier-Domain Analysis of BHEX for Photon-Ring Inference}, 2026.

\noindent{}\hypertarget{ref-2}{}{[}2{]} Jonathan R. Landers,
\emph{Geodesic-Provenance Coherence Bounds for Fourier-Domain Photon-Ring Inference}, 2026.

\noindent{}\hypertarget{ref-3}{}{[}3{]} Black Hole Explorer Collaboration,
\emph{Black Hole Explorer: Motivation and Vision}, arXiv:2406.12917.
\url{https://arxiv.org/abs/2406.12917}

\noindent{}\hypertarget{ref-4}{}{[}4{]} Black Hole Explorer Collaboration,
\emph{The Black Hole Explorer: Photon Ring Science, Detection and Shape Measurement},
arXiv:2406.09498. \url{https://arxiv.org/abs/2406.09498}

\noindent{}\hypertarget{ref-5}{}{[}5{]} Black Hole Explorer Collaboration,
\emph{Interferometric Inference of Black Hole Spin from Photon Ring Size and Brightness},
arXiv:2509.23628. \url{https://arxiv.org/abs/2509.23628}

\end{document}
